% Context from chapters/denoising.tex; for mathematical reference, not compilation.
\section{Signal-Dependent Noise}
\label{sect-data-dept-noises}

In many imaging systems, the variance of the noise affecting $f_{0,n}$ depends on the intensity $f_{0,n}$ itself. We now study Poisson counts and multiplicative noise. Both can be expressed as additive residuals, but their residual distributions depend on the signal.

                                                     
\subsection{Poisson Noise}

Many imaging devices measure intensity by counting photons. Examples include digital cameras, confocal microscopy, PET and SPECT tomography.

  
\paragraph{Poisson model.}

For an unknown mean intensity $f_{0,n}\geq0$, a Poisson model describes the observed photon count:
\eq{
	f_n \sim \Pp(\la)
	\qwhereq \lambda=f_{0,n}\in[0,\infty),
}
The Poisson distribution $\Pp(\la)$ has probability mass function
\eq{
	\PP(f_n=k)=\frac{\la^k e^{-\la}}{k!}
}
and its parameter varies across pixels with the underlying intensity. Figure \ref{fig-poisson-distrib} shows several Poisson distributions.

\myfigure{
\image{denoising}{.5}{poisson-distrib}
}{ 
	 Poisson distributions for various $\la$.   
}{fig-poisson-distrib}

One has 
\eq{
	\mathbb E(f_n)=\la=f_{0,n}
	\qandq
	\text{Var}(f_n)=\la=f_{0,n}
}
Thus denoising again estimates a mean from one observation, but the variance now varies with intensity. The absolute noise level increases as more photons reach the sensor. For positive intensity, the relative standard deviation $f_{0,n}^{-1/2}$ instead decreases.

Figure \ref{fig-poisson-denoising} shows Poisson-corrupted versions of a clean image $f_0$ at different maximum mean intensities $\la_{\max}$.

\myfigure{
\tabquatre{
\image{denoising}{.23}{poisson-denoising-1}&
\image{denoising}{.23}{poisson-denoising-2}&
\image{denoising}{.23}{poisson-denoising-3}&
\image{denoising}{.23}{poisson-denoising-4}\\
Example 1 & Example 2 & Example 3 & Example 4
}
}{ 
	 Examples of Poisson-corrupted images at different count levels. Numerical acquisition parameters are not available for these legacy illustrations.  
}{fig-poisson-denoising}

  
\paragraph{Variance stabilization.}

Applying a thresholding estimator
\eq{
	\Dd(f) = \sum_m S_T^q(\dotp{f}{\psi_m}) \psi_m
}
directly to $f$ can perform poorly because a single threshold $T$ cannot adapt to spatially varying noise. A variance-stabilizing transform $\phi : [0,+\infty) \rightarrow \RR$, applied pixelwise, maps the counts to $\phi(f)$. An additive Gaussian white-noise model can then be a useful approximation:
\eql{\label{eq-var-stab}
	\phi(f) \approx \phi(f_0) + w
}
where $w_n\sim\Nn(0,\sigma^2)$ are independent centered Gaussians with constant variance $\si^2$.

The model~\eqref{eq-var-stab} is approximate, and its accuracy depends on the intensity range of $f_{0,n}$.
Two common variance-stabilizing transforms for Poisson noise are the Anscombe transform
\eq{
	\phi(x) = 2\sqrt{x+3/8}
}
and the Freeman--Tukey transform
\eq{
	\phi(x) = \sqrt{x+1}+\sqrt{x}.
}
Figure \ref{fig-variance-stabilization} compares the variance of $\phi(f)$ after these transformations.

\myfigure{
\image{denoising}{.5}{variance-stabilization-corrected}
}{ 
	 Exact Poisson variances of the Anscombe, Freeman--Tukey, and $2\sqrt{x}$ transforms as functions of the mean intensity.  
}{fig-variance-stabilization}

A variance-stabilized denoiser is defined as
\eq{
	\Dd^{\text{stab},q}(f) =  \phi^{-1}( \sum_m S_T^q(\dotp{\phi(f)}{\psi_m}) \psi_m )
}
where $\phi^{-1}$ is applied after projecting the denoised values onto its domain. The direct inverse generally introduces bias, especially at low counts; bias-corrected inverse transforms can improve the result.

At the moderate count levels in Figure \ref{fig-poisson-wav}, variance stabilization improves the denoising result.

\myfigure{
\tabtrois{
\image{denoising}{.32}{poisson-wav-noisy}&
\image{denoising}{.32}{poisson-wav-unstab}&
\image{denoising}{.32}{poisson-wav-stab}
}
}{ 
	Left: noisy image, center: denoising without variance stabilization, right: denoising after variance stabilization.   
}{fig-poisson-wav}

                